\documentclass[11pt,letterpaper]{article}
\usepackage[technical-reference]{cornell-notes}
\usepackage{needspace}

\title{Clause 3: Terms and Definitions - Cornell Notes}
\author{}
\date{}
\setDocTitle{Clause 3: Terms and Definitions}
\setDocSubtitle{C++ 2024 Cornell Notes}
\setDocOwner{C++ 2024 Cornell Notes}
\setDocAuthor{}
\setDocDate{}
\setCornellCollection{C++ 2024 Cornell Notes}
\setCornellUnitType{Clause}
\setCornellUnitNumber{3}
\setCornellUnitTitle{Terms and Definitions}
\setCornellCompanionLabel{C++ Technical Reference}
\hypersetup{pdftitle={Clause 3: Terms and Definitions - Cornell Notes},pdfsubject={C++ technical study notes},pdfauthor={}}

\begin{document}
\maketitle
\begin{CornellReferenceOrientationBox}
\textbf{Purpose.} Defines cross-cutting vocabulary used by the core language and library. The definitions distinguish program status and behavior categories, execution and concurrency concepts, library roles, regular-expression terms, and the components of signatures.
\par\textbf{Role.} Provides the shared semantic vocabulary needed to read requirements consistently.
\par\textbf{Principal dependencies.} External vocabulary and mathematical terminology are supplemented by definitions specific to this source material.
\par\textbf{Primary audience.} all readers, especially reviewers classifying behavior or library contracts.
\par\textbf{Major subject groups.} access; argument; argument; argument; argument; block; block; C standard library; character; character container type; collating element; component; and related subclauses.
\end{CornellReferenceOrientationBox}
\section{Learning Objectives}
\begin{itemize}[label=\textcolor{CornellReferenceTeal}{$\blacktriangleright$}]
\item Define the governing ideas of access and state the consequences for a program or implementation.
\item Identify the governing ideas of argument and state the consequences for a program or implementation.
\item Distinguish the governing ideas of argument and state the consequences for a program or implementation.
\item Explain the governing ideas of argument and state the consequences for a program or implementation.
\item Trace the governing ideas of argument and state the consequences for a program or implementation.
\item Apply the governing ideas of block and state the consequences for a program or implementation.
\item Analyze the governing ideas of block and state the consequences for a program or implementation.
\item Evaluate the governing ideas of C standard library and state the consequences for a program or implementation.
\item Diagnose the governing ideas of character and state the consequences for a program or implementation.
\item Compare the governing ideas of character container type and state the consequences for a program or implementation.
\item Verify the governing ideas of collating element and state the consequences for a program or implementation.
\item Relate the governing ideas of component and state the consequences for a program or implementation.
\item Classify the governing ideas of conditionally-supported and state the consequences for a program or implementation.
\item Justify the governing ideas of constant subexpression and state the consequences for a program or implementation.
\item Audit the governing ideas of deadlock and state the consequences for a program or implementation.
\item Summarize the governing ideas of default behavior and state the consequences for a program or implementation.
\end{itemize}
\section{Normative Structure Map}
\small\begin{longtable}{@{}>{\RaggedRight\arraybackslash}p{0.13\textwidth}>{\RaggedRight\arraybackslash}p{0.27\textwidth}>{\RaggedRight\arraybackslash}p{0.19\textwidth}>{\RaggedRight\arraybackslash}p{0.31\textwidth}@{}}
\toprule\textbf{Subclause} & \textbf{Subject} & \textbf{Rule category} & \textbf{Key dependencies / entities}\\\midrule\endfirsthead
\toprule\textbf{Subclause} & \textbf{Subject} & \textbf{Rule category} & \textbf{Key dependencies / entities}\\\midrule\endhead
\bottomrule\endfoot
3.1 & access & core-language semantics & External vocabulary and mathematical terminology are supplemented by definitions specific to this source material.\\
3.2 & argument & core-language semantics & External vocabulary and mathematical terminology are supplemented by definitions specific to this source material.\\
3.3 & argument & core-language semantics & External vocabulary and mathematical terminology are supplemented by definitions specific to this source material.\\
3.4 & argument & core-language semantics & External vocabulary and mathematical terminology are supplemented by definitions specific to this source material.\\
3.5 & argument & core-language semantics & External vocabulary and mathematical terminology are supplemented by definitions specific to this source material.\\
3.6 & block & concurrency contract & External vocabulary and mathematical terminology are supplemented by definitions specific to this source material.\\
3.7 & block & concurrency contract & External vocabulary and mathematical terminology are supplemented by definitions specific to this source material.\\
3.8 & C standard library & core-language semantics & External vocabulary and mathematical terminology are supplemented by definitions specific to this source material.\\
3.9 & character & core-language semantics & External vocabulary and mathematical terminology are supplemented by definitions specific to this source material.\\
3.10 & character container type & core-language semantics & External vocabulary and mathematical terminology are supplemented by definitions specific to this source material.\\
3.11 & collating element & core-language semantics & External vocabulary and mathematical terminology are supplemented by definitions specific to this source material.\\
3.12 & component & core-language semantics & External vocabulary and mathematical terminology are supplemented by definitions specific to this source material.\\
3.13 & conditionally-supported & concurrency contract & External vocabulary and mathematical terminology are supplemented by definitions specific to this source material.\\
3.14 & constant subexpression & core-language semantics & External vocabulary and mathematical terminology are supplemented by definitions specific to this source material.\\
3.15 & deadlock & concurrency contract & External vocabulary and mathematical terminology are supplemented by definitions specific to this source material.\\
3.16 & default behavior & core-language semantics & External vocabulary and mathematical terminology are supplemented by definitions specific to this source material.\\
3.17 & diagnostic message & error behavior & External vocabulary and mathematical terminology are supplemented by definitions specific to this source material.\\
3.18 & dynamic type & core-language semantics & External vocabulary and mathematical terminology are supplemented by definitions specific to this source material.\\
3.19 & dynamic type & core-language semantics & External vocabulary and mathematical terminology are supplemented by definitions specific to this source material.\\
3.20 & expression-equivalent & core-language semantics & External vocabulary and mathematical terminology are supplemented by definitions specific to this source material.\\
3.21 & finite state machine & core-language semantics & External vocabulary and mathematical terminology are supplemented by definitions specific to this source material.\\
3.22 & format specifier & core-language semantics & External vocabulary and mathematical terminology are supplemented by definitions specific to this source material.\\
3.23 & handler function & operation semantics & External vocabulary and mathematical terminology are supplemented by definitions specific to this source material.\\
3.24 & ill-formed program & core-language semantics & External vocabulary and mathematical terminology are supplemented by definitions specific to this source material.\\
3.25 & implementation-defined behavior & core-language semantics & External vocabulary and mathematical terminology are supplemented by definitions specific to this source material.\\
3.26 & implementation-defined strict total order over pointers & core-language semantics & External vocabulary and mathematical terminology are supplemented by definitions specific to this source material.\\
3.27 & implementation limit & core-language semantics & External vocabulary and mathematical terminology are supplemented by definitions specific to this source material.\\
3.28 & locale-specific behavior & core-language semantics & External vocabulary and mathematical terminology are supplemented by definitions specific to this source material.\\
3.29 & matched & core-language semantics & External vocabulary and mathematical terminology are supplemented by definitions specific to this source material.\\
3.30 & modifier function & operation semantics & External vocabulary and mathematical terminology are supplemented by definitions specific to this source material.\\
3.31 & move assignment & object contract & External vocabulary and mathematical terminology are supplemented by definitions specific to this source material.\\
3.32 & move construction & core-language semantics & External vocabulary and mathematical terminology are supplemented by definitions specific to this source material.\\
3.33 & non-constant library call & core-language semantics & External vocabulary and mathematical terminology are supplemented by definitions specific to this source material.\\
3.34 & NTCTS & core-language semantics & External vocabulary and mathematical terminology are supplemented by definitions specific to this source material.\\
3.35 & observer function & operation semantics & External vocabulary and mathematical terminology are supplemented by definitions specific to this source material.\\
3.36 & parameter & core-language semantics & External vocabulary and mathematical terminology are supplemented by definitions specific to this source material.\\
3.37 & parameter & core-language semantics & External vocabulary and mathematical terminology are supplemented by definitions specific to this source material.\\
3.38 & parameter & core-language semantics & External vocabulary and mathematical terminology are supplemented by definitions specific to this source material.\\
3.39 & primary equivalence class & core-language semantics & External vocabulary and mathematical terminology are supplemented by definitions specific to this source material.\\
3.40 & program-defined specialization & core-language semantics & External vocabulary and mathematical terminology are supplemented by definitions specific to this source material.\\
3.41 & program-defined type & core-language semantics & External vocabulary and mathematical terminology are supplemented by definitions specific to this source material.\\
3.42 & projection & core-language semantics & External vocabulary and mathematical terminology are supplemented by definitions specific to this source material.\\
3.43 & referenceable type & core-language semantics & External vocabulary and mathematical terminology are supplemented by definitions specific to this source material.\\
3.44 & regular expression & core-language semantics & External vocabulary and mathematical terminology are supplemented by definitions specific to this source material.\\
3.45 & replacement function & operation semantics & External vocabulary and mathematical terminology are supplemented by definitions specific to this source material.\\
3.46 & required behavior & core-language semantics & External vocabulary and mathematical terminology are supplemented by definitions specific to this source material.\\
3.47 & reserved function & operation semantics & External vocabulary and mathematical terminology are supplemented by definitions specific to this source material.\\
3.48 & signature & core-language semantics & External vocabulary and mathematical terminology are supplemented by definitions specific to this source material.\\
3.49 & signature & core-language semantics & External vocabulary and mathematical terminology are supplemented by definitions specific to this source material.\\
3.50 & signature & core-language semantics & External vocabulary and mathematical terminology are supplemented by definitions specific to this source material.\\
3.51 & signature & core-language semantics & External vocabulary and mathematical terminology are supplemented by definitions specific to this source material.\\
3.52 & signature & core-language semantics & External vocabulary and mathematical terminology are supplemented by definitions specific to this source material.\\
3.53 & signature & core-language semantics & External vocabulary and mathematical terminology are supplemented by definitions specific to this source material.\\
3.54 & signature & core-language semantics & External vocabulary and mathematical terminology are supplemented by definitions specific to this source material.\\
3.55 & signature & core-language semantics & External vocabulary and mathematical terminology are supplemented by definitions specific to this source material.\\
3.56 & signature & core-language semantics & External vocabulary and mathematical terminology are supplemented by definitions specific to this source material.\\
3.57 & stable algorithm & operation semantics & External vocabulary and mathematical terminology are supplemented by definitions specific to this source material.\\
3.58 & static type & core-language semantics & External vocabulary and mathematical terminology are supplemented by definitions specific to this source material.\\
3.59 & sub-expression & core-language semantics & External vocabulary and mathematical terminology are supplemented by definitions specific to this source material.\\
3.60 & traits class & core-language semantics & External vocabulary and mathematical terminology are supplemented by definitions specific to this source material.\\
3.61 & unblock & concurrency contract & External vocabulary and mathematical terminology are supplemented by definitions specific to this source material.\\
3.62 & undefined behavior & core-language semantics & External vocabulary and mathematical terminology are supplemented by definitions specific to this source material.\\
3.63 & unspecified behavior & core-language semantics & External vocabulary and mathematical terminology are supplemented by definitions specific to this source material.\\
3.64 & valid but unspecified state & core-language semantics & External vocabulary and mathematical terminology are supplemented by definitions specific to this source material.\\
3.65 & well-formed program & core-language semantics & External vocabulary and mathematical terminology are supplemented by definitions specific to this source material.\\
\end{longtable}\normalsize
\begin{CornellReferenceNormativeBox}A heading maps the clause; it does not replace the detailed conditions. For each operation, read requirements, effects, result, exceptions, complexity, remarks, and notes with their stated normative force.\end{CornellReferenceNormativeBox}
\subsection*{Rule Classification Guide}
\small\begin{longtable}{@{}>{\RaggedRight\arraybackslash}p{0.25\textwidth}>{\RaggedRight\arraybackslash}p{0.67\textwidth}@{}}\toprule\textbf{Label or category} & \textbf{How to read it}\\\midrule\endfirsthead\toprule\textbf{Label or category} & \textbf{How to read it (continued)}\\\midrule\endhead\bottomrule\endfoot
Well-formed / ill-formed & Formation is governed by syntax and semantic rules; an ill-formed program does not become well-formed because one implementation accepts it.\\
Diagnosable rule & A violation requires at least one diagnostic unless the wording places the case outside the diagnosable set.\\
No diagnostic required & The program can be ill-formed while the implementation has no diagnostic obligation.\\
Undefined behavior & No execution requirements are imposed; translation success does not establish safety or portability.\\
Unspecified behavior & One of the permitted outcomes may be chosen without documentation.\\
Implementation-defined behavior & The implementation chooses and documents the behavior.\\
Conditionally supported & The implementation may omit support but documents unsupported constructs.\\
\end{longtable}\normalsize
\section{Cornell Cue-and-Notes}
\Needspace{8\baselineskip}
\subsection*{3.1 access}
\begin{longtable}{@{}>{\RaggedRight\arraybackslash\columncolor{CornellReferenceCueGray}}p{0.28\textwidth}>{\RaggedRight\arraybackslash}p{0.64\textwidth}@{}}
\rowcolor{CornellReferenceNavy}\color{white}\textbf{Cue / recall prompt} & \color{white}\textbf{Notes}\\\endfirsthead
\rowcolor{CornellReferenceNavy}\color{white}\textbf{Cue / recall prompt} & \color{white}\textbf{Notes (continued)}\\\endhead
\arrayrulecolor{CornellReferenceLineGray}
\textbf{What rule applies to access?}\\[-1mm]{\scriptsize\CornellClauseRef{3.1} [defns.\allowbreak{}access]} & An execution-time access is a read or modification of an object's value; class operations can contain accesses without the call itself automatically being one.\\\hline
\end{longtable}
\begin{CornellReferenceSummaryBox}access supplies the core-language semantics portion of this clause. The key review move is to connect its local wording to the clause-wide model, then classify each condition by normative force before relying on it.\end{CornellReferenceSummaryBox}
\Needspace{8\baselineskip}
\subsection*{3.2 argument}
\begin{longtable}{@{}>{\RaggedRight\arraybackslash\columncolor{CornellReferenceCueGray}}p{0.28\textwidth}>{\RaggedRight\arraybackslash}p{0.64\textwidth}@{}}
\rowcolor{CornellReferenceNavy}\color{white}\textbf{Cue / recall prompt} & \color{white}\textbf{Notes}\\\endfirsthead
\rowcolor{CornellReferenceNavy}\color{white}\textbf{Cue / recall prompt} & \color{white}\textbf{Notes (continued)}\\\endhead
\arrayrulecolor{CornellReferenceLineGray}
\textbf{What rule applies to argument?}\\[-1mm]{\scriptsize\CornellClauseRef{3.2} [defns.\allowbreak{}argument]} & In a function call, an argument is an expression supplied between the call parentheses.\\\hline
\end{longtable}
\begin{CornellReferenceSummaryBox}argument supplies the core-language semantics portion of this clause. The key review move is to connect its local wording to the clause-wide model, then classify each condition by normative force before relying on it.\end{CornellReferenceSummaryBox}
\Needspace{8\baselineskip}
\subsection*{3.3 argument}
\begin{longtable}{@{}>{\RaggedRight\arraybackslash\columncolor{CornellReferenceCueGray}}p{0.28\textwidth}>{\RaggedRight\arraybackslash}p{0.64\textwidth}@{}}
\rowcolor{CornellReferenceNavy}\color{white}\textbf{Cue / recall prompt} & \color{white}\textbf{Notes}\\\endfirsthead
\rowcolor{CornellReferenceNavy}\color{white}\textbf{Cue / recall prompt} & \color{white}\textbf{Notes (continued)}\\\endhead
\arrayrulecolor{CornellReferenceLineGray}
\textbf{What rule applies to argument?}\\[-1mm]{\scriptsize\CornellClauseRef{3.3} [defns.\allowbreak{}argument.\allowbreak{}macro]} & For a function-like macro, an argument is the preprocessing-token sequence supplied for one macro parameter.\\\hline
\end{longtable}
\begin{CornellReferenceSummaryBox}argument supplies the core-language semantics portion of this clause. The key review move is to connect its local wording to the clause-wide model, then classify each condition by normative force before relying on it.\end{CornellReferenceSummaryBox}
\Needspace{8\baselineskip}
\subsection*{3.4 argument}
\begin{longtable}{@{}>{\RaggedRight\arraybackslash\columncolor{CornellReferenceCueGray}}p{0.28\textwidth}>{\RaggedRight\arraybackslash}p{0.64\textwidth}@{}}
\rowcolor{CornellReferenceNavy}\color{white}\textbf{Cue / recall prompt} & \color{white}\textbf{Notes}\\\endfirsthead
\rowcolor{CornellReferenceNavy}\color{white}\textbf{Cue / recall prompt} & \color{white}\textbf{Notes (continued)}\\\endhead
\arrayrulecolor{CornellReferenceLineGray}
\textbf{What rule applies to argument?}\\[-1mm]{\scriptsize\CornellClauseRef{3.4} [defns.\allowbreak{}argument.\allowbreak{}throw]} & For a throw-expression, the argument is the operand used to initialize the exception object.\\\hline
\end{longtable}
\begin{CornellReferenceSummaryBox}argument supplies the core-language semantics portion of this clause. The key review move is to connect its local wording to the clause-wide model, then classify each condition by normative force before relying on it.\end{CornellReferenceSummaryBox}
\Needspace{8\baselineskip}
\subsection*{3.5 argument}
\begin{longtable}{@{}>{\RaggedRight\arraybackslash\columncolor{CornellReferenceCueGray}}p{0.28\textwidth}>{\RaggedRight\arraybackslash}p{0.64\textwidth}@{}}
\rowcolor{CornellReferenceNavy}\color{white}\textbf{Cue / recall prompt} & \color{white}\textbf{Notes}\\\endfirsthead
\rowcolor{CornellReferenceNavy}\color{white}\textbf{Cue / recall prompt} & \color{white}\textbf{Notes (continued)}\\\endhead
\arrayrulecolor{CornellReferenceLineGray}
\textbf{What rule applies to argument?}\\[-1mm]{\scriptsize\CornellClauseRef{3.5} [defns.\allowbreak{}argument.\allowbreak{}templ]} & For template instantiation, an argument is the type, constant expression, or qualifying name supplied to a template parameter.\\\hline
\end{longtable}
\begin{CornellReferenceSummaryBox}argument supplies the core-language semantics portion of this clause. The key review move is to connect its local wording to the clause-wide model, then classify each condition by normative force before relying on it.\end{CornellReferenceSummaryBox}
\Needspace{8\baselineskip}
\subsection*{3.6 block}
\begin{longtable}{@{}>{\RaggedRight\arraybackslash\columncolor{CornellReferenceCueGray}}p{0.28\textwidth}>{\RaggedRight\arraybackslash}p{0.64\textwidth}@{}}
\rowcolor{CornellReferenceNavy}\color{white}\textbf{Cue / recall prompt} & \color{white}\textbf{Notes}\\\endfirsthead
\rowcolor{CornellReferenceNavy}\color{white}\textbf{Cue / recall prompt} & \color{white}\textbf{Notes (continued)}\\\endhead
\arrayrulecolor{CornellReferenceLineGray}
\textbf{What rule applies to block?}\\[-1mm]{\scriptsize\CornellClauseRef{3.6} [defns.\allowbreak{}block]} & In execution, blocking means waiting for an external condition rather than merely waiting for the implementation to schedule ordinary execution steps.\\\hline
\end{longtable}
\begin{CornellReferenceSummaryBox}block supplies the concurrency contract portion of this clause. The key review move is to connect its local wording to the clause-wide model, then classify each condition by normative force before relying on it.\end{CornellReferenceSummaryBox}
\Needspace{8\baselineskip}
\subsection*{3.7 block}
\begin{longtable}{@{}>{\RaggedRight\arraybackslash\columncolor{CornellReferenceCueGray}}p{0.28\textwidth}>{\RaggedRight\arraybackslash}p{0.64\textwidth}@{}}
\rowcolor{CornellReferenceNavy}\color{white}\textbf{Cue / recall prompt} & \color{white}\textbf{Notes}\\\endfirsthead
\rowcolor{CornellReferenceNavy}\color{white}\textbf{Cue / recall prompt} & \color{white}\textbf{Notes (continued)}\\\endhead
\arrayrulecolor{CornellReferenceLineGray}
\textbf{What rule applies to block?}\\[-1mm]{\scriptsize\CornellClauseRef{3.7} [defns.\allowbreak{}block.\allowbreak{}stmt]} & In statement grammar, a block is a compound statement.\\\hline
\end{longtable}
\begin{CornellReferenceSummaryBox}block supplies the concurrency contract portion of this clause. The key review move is to connect its local wording to the clause-wide model, then classify each condition by normative force before relying on it.\end{CornellReferenceSummaryBox}
\Needspace{8\baselineskip}
\subsection*{3.8 C standard library}
\begin{longtable}{@{}>{\RaggedRight\arraybackslash\columncolor{CornellReferenceCueGray}}p{0.28\textwidth}>{\RaggedRight\arraybackslash}p{0.64\textwidth}@{}}
\rowcolor{CornellReferenceNavy}\color{white}\textbf{Cue / recall prompt} & \color{white}\textbf{Notes}\\\endfirsthead
\rowcolor{CornellReferenceNavy}\color{white}\textbf{Cue / recall prompt} & \color{white}\textbf{Notes (continued)}\\\endhead
\arrayrulecolor{CornellReferenceLineGray}
\textbf{What rule applies to C standard library?}\\[-1mm]{\scriptsize\CornellClauseRef{3.8} [defns.\allowbreak{}c.\allowbreak{}lib]} & The inherited C library forms a qualified subset of the C++ library; later library clauses state the adaptations.\\\hline
\end{longtable}
\begin{CornellReferenceSummaryBox}C standard library supplies the core-language semantics portion of this clause. The key review move is to connect its local wording to the clause-wide model, then classify each condition by normative force before relying on it.\end{CornellReferenceSummaryBox}
\Needspace{8\baselineskip}
\subsection*{3.9 character}
\begin{longtable}{@{}>{\RaggedRight\arraybackslash\columncolor{CornellReferenceCueGray}}p{0.28\textwidth}>{\RaggedRight\arraybackslash}p{0.64\textwidth}@{}}
\rowcolor{CornellReferenceNavy}\color{white}\textbf{Cue / recall prompt} & \color{white}\textbf{Notes}\\\endfirsthead
\rowcolor{CornellReferenceNavy}\color{white}\textbf{Cue / recall prompt} & \color{white}\textbf{Notes (continued)}\\\endhead
\arrayrulecolor{CornellReferenceLineGray}
\textbf{What rule applies to character?}\\[-1mm]{\scriptsize\CornellClauseRef{3.9} [defns.\allowbreak{}character]} & A library character is a value used in a sequence to represent text and is not limited to the built-in character types.\\\hline
\end{longtable}
\begin{CornellReferenceSummaryBox}character supplies the core-language semantics portion of this clause. The key review move is to connect its local wording to the clause-wide model, then classify each condition by normative force before relying on it.\end{CornellReferenceSummaryBox}
\Needspace{8\baselineskip}
\subsection*{3.10 character container type}
\begin{longtable}{@{}>{\RaggedRight\arraybackslash\columncolor{CornellReferenceCueGray}}p{0.28\textwidth}>{\RaggedRight\arraybackslash}p{0.64\textwidth}@{}}
\rowcolor{CornellReferenceNavy}\color{white}\textbf{Cue / recall prompt} & \color{white}\textbf{Notes}\\\endfirsthead
\rowcolor{CornellReferenceNavy}\color{white}\textbf{Cue / recall prompt} & \color{white}\textbf{Notes (continued)}\\\endhead
\arrayrulecolor{CornellReferenceLineGray}
\textbf{What rule applies to character container type?}\\[-1mm]{\scriptsize\CornellClauseRef{3.10} [defns.\allowbreak{}character.\allowbreak{}container]} & A character container type is a class or type used as the representation type for characters in text-related templates.\\\hline
\end{longtable}
\begin{CornellReferenceSummaryBox}character container type supplies the core-language semantics portion of this clause. The key review move is to connect its local wording to the clause-wide model, then classify each condition by normative force before relying on it.\end{CornellReferenceSummaryBox}
\Needspace{8\baselineskip}
\subsection*{3.11 collating element}
\begin{longtable}{@{}>{\RaggedRight\arraybackslash\columncolor{CornellReferenceCueGray}}p{0.28\textwidth}>{\RaggedRight\arraybackslash}p{0.64\textwidth}@{}}
\rowcolor{CornellReferenceNavy}\color{white}\textbf{Cue / recall prompt} & \color{white}\textbf{Notes}\\\endfirsthead
\rowcolor{CornellReferenceNavy}\color{white}\textbf{Cue / recall prompt} & \color{white}\textbf{Notes (continued)}\\\endhead
\arrayrulecolor{CornellReferenceLineGray}
\textbf{What rule applies to collating element?}\\[-1mm]{\scriptsize\CornellClauseRef{3.11} [defns.\allowbreak{}regex.\allowbreak{}collating.\allowbreak{}element]} & A collating element is one or more characters treated as a single ordering unit by the active locale.\\\hline
\end{longtable}
\begin{CornellReferenceSummaryBox}collating element supplies the core-language semantics portion of this clause. The key review move is to connect its local wording to the clause-wide model, then classify each condition by normative force before relying on it.\end{CornellReferenceSummaryBox}
\Needspace{8\baselineskip}
\subsection*{3.12 component}
\begin{longtable}{@{}>{\RaggedRight\arraybackslash\columncolor{CornellReferenceCueGray}}p{0.28\textwidth}>{\RaggedRight\arraybackslash}p{0.64\textwidth}@{}}
\rowcolor{CornellReferenceNavy}\color{white}\textbf{Cue / recall prompt} & \color{white}\textbf{Notes}\\\endfirsthead
\rowcolor{CornellReferenceNavy}\color{white}\textbf{Cue / recall prompt} & \color{white}\textbf{Notes (continued)}\\\endhead
\arrayrulecolor{CornellReferenceLineGray}
\textbf{What rule applies to component?}\\[-1mm]{\scriptsize\CornellClauseRef{3.12} [defns.\allowbreak{}component]} & A library component is a directly related family of entities connected through membership, parameters, or return types.\\\hline
\end{longtable}
\begin{CornellReferenceSummaryBox}component supplies the core-language semantics portion of this clause. The key review move is to connect its local wording to the clause-wide model, then classify each condition by normative force before relying on it.\end{CornellReferenceSummaryBox}
\Needspace{8\baselineskip}
\subsection*{3.13 conditionally-supported}
\begin{longtable}{@{}>{\RaggedRight\arraybackslash\columncolor{CornellReferenceCueGray}}p{0.28\textwidth}>{\RaggedRight\arraybackslash}p{0.64\textwidth}@{}}
\rowcolor{CornellReferenceNavy}\color{white}\textbf{Cue / recall prompt} & \color{white}\textbf{Notes}\\\endfirsthead
\rowcolor{CornellReferenceNavy}\color{white}\textbf{Cue / recall prompt} & \color{white}\textbf{Notes (continued)}\\\endhead
\arrayrulecolor{CornellReferenceLineGray}
\textbf{What rule applies to conditionally-supported?}\\[-1mm]{\scriptsize\CornellClauseRef{3.13} [defns.\allowbreak{}cond.\allowbreak{}supp]} & A conditionally supported construct need not be provided, but unsupported constructs in this category must be documented by the implementation.\\\hline
\end{longtable}
\begin{CornellReferenceSummaryBox}conditionally-supported supplies the concurrency contract portion of this clause. The key review move is to connect its local wording to the clause-wide model, then classify each condition by normative force before relying on it.\end{CornellReferenceSummaryBox}
\Needspace{8\baselineskip}
\subsection*{3.14 constant subexpression}
\begin{longtable}{@{}>{\RaggedRight\arraybackslash\columncolor{CornellReferenceCueGray}}p{0.28\textwidth}>{\RaggedRight\arraybackslash}p{0.64\textwidth}@{}}
\rowcolor{CornellReferenceNavy}\color{white}\textbf{Cue / recall prompt} & \color{white}\textbf{Notes}\\\endfirsthead
\rowcolor{CornellReferenceNavy}\color{white}\textbf{Cue / recall prompt} & \color{white}\textbf{Notes (continued)}\\\endhead
\arrayrulecolor{CornellReferenceLineGray}
\textbf{What rule applies to constant subexpression?}\\[-1mm]{\scriptsize\CornellClauseRef{3.14} [defns.\allowbreak{}const.\allowbreak{}subexpr]} & A constant subexpression can be evaluated within a conditional expression without preventing that enclosing expression from being a core constant expression.\\\hline
\end{longtable}
\begin{CornellReferenceSummaryBox}constant subexpression supplies the core-language semantics portion of this clause. The key review move is to connect its local wording to the clause-wide model, then classify each condition by normative force before relying on it.\end{CornellReferenceSummaryBox}
\Needspace{8\baselineskip}
\subsection*{3.15 deadlock}
\begin{longtable}{@{}>{\RaggedRight\arraybackslash\columncolor{CornellReferenceCueGray}}p{0.28\textwidth}>{\RaggedRight\arraybackslash}p{0.64\textwidth}@{}}
\rowcolor{CornellReferenceNavy}\color{white}\textbf{Cue / recall prompt} & \color{white}\textbf{Notes}\\\endfirsthead
\rowcolor{CornellReferenceNavy}\color{white}\textbf{Cue / recall prompt} & \color{white}\textbf{Notes (continued)}\\\endhead
\arrayrulecolor{CornellReferenceLineGray}
\textbf{What rule applies to deadlock?}\\[-1mm]{\scriptsize\CornellClauseRef{3.15} [defns.\allowbreak{}deadlock]} & Deadlock occurs when threads cannot proceed because each waits for a condition that depends on another blocked participant.\\\hline
\end{longtable}
\begin{CornellReferenceSummaryBox}deadlock supplies the concurrency contract portion of this clause. The key review move is to connect its local wording to the clause-wide model, then classify each condition by normative force before relying on it.\end{CornellReferenceSummaryBox}
\Needspace{8\baselineskip}
\subsection*{3.16 default behavior}
\begin{longtable}{@{}>{\RaggedRight\arraybackslash\columncolor{CornellReferenceCueGray}}p{0.28\textwidth}>{\RaggedRight\arraybackslash}p{0.64\textwidth}@{}}
\rowcolor{CornellReferenceNavy}\color{white}\textbf{Cue / recall prompt} & \color{white}\textbf{Notes}\\\endfirsthead
\rowcolor{CornellReferenceNavy}\color{white}\textbf{Cue / recall prompt} & \color{white}\textbf{Notes (continued)}\\\endhead
\arrayrulecolor{CornellReferenceLineGray}
\textbf{What rule applies to default behavior?}\\[-1mm]{\scriptsize\CornellClauseRef{3.16} [defns.\allowbreak{}default.\allowbreak{}behavior.\allowbreak{}impl]} & Default behavior is the implementation's chosen behavior within the bounds of the required behavior for a replaceable or installable facility.\\\hline
\end{longtable}
\begin{CornellReferenceSummaryBox}default behavior supplies the core-language semantics portion of this clause. The key review move is to connect its local wording to the clause-wide model, then classify each condition by normative force before relying on it.\end{CornellReferenceSummaryBox}
\Needspace{8\baselineskip}
\subsection*{3.17 diagnostic message}
\begin{longtable}{@{}>{\RaggedRight\arraybackslash\columncolor{CornellReferenceCueGray}}p{0.28\textwidth}>{\RaggedRight\arraybackslash}p{0.64\textwidth}@{}}
\rowcolor{CornellReferenceNavy}\color{white}\textbf{Cue / recall prompt} & \color{white}\textbf{Notes}\\\endfirsthead
\rowcolor{CornellReferenceNavy}\color{white}\textbf{Cue / recall prompt} & \color{white}\textbf{Notes (continued)}\\\endhead
\arrayrulecolor{CornellReferenceLineGray}
\textbf{What rule applies to diagnostic message?}\\[-1mm]{\scriptsize\CornellClauseRef{3.17} [defns.\allowbreak{}diagnostic]} & A diagnostic message is any message belonging to the implementation-defined subset designated as diagnostics.\\\hline
\end{longtable}
\begin{CornellReferenceSummaryBox}diagnostic message supplies the error behavior portion of this clause. The key review move is to connect its local wording to the clause-wide model, then classify each condition by normative force before relying on it.\end{CornellReferenceSummaryBox}
\Needspace{8\baselineskip}
\subsection*{3.18 dynamic type}
\begin{longtable}{@{}>{\RaggedRight\arraybackslash\columncolor{CornellReferenceCueGray}}p{0.28\textwidth}>{\RaggedRight\arraybackslash}p{0.64\textwidth}@{}}
\rowcolor{CornellReferenceNavy}\color{white}\textbf{Cue / recall prompt} & \color{white}\textbf{Notes}\\\endfirsthead
\rowcolor{CornellReferenceNavy}\color{white}\textbf{Cue / recall prompt} & \color{white}\textbf{Notes (continued)}\\\endhead
\arrayrulecolor{CornellReferenceLineGray}
\textbf{What rule applies to dynamic type?}\\[-1mm]{\scriptsize\CornellClauseRef{3.18} [defns.\allowbreak{}dynamic.\allowbreak{}type]} & For a glvalue, dynamic type is the type of the most-derived object denoted at runtime.\\\hline
\end{longtable}
\begin{CornellReferenceSummaryBox}dynamic type supplies the core-language semantics portion of this clause. The key review move is to connect its local wording to the clause-wide model, then classify each condition by normative force before relying on it.\end{CornellReferenceSummaryBox}
\Needspace{8\baselineskip}
\subsection*{3.19 dynamic type}
\begin{longtable}{@{}>{\RaggedRight\arraybackslash\columncolor{CornellReferenceCueGray}}p{0.28\textwidth}>{\RaggedRight\arraybackslash}p{0.64\textwidth}@{}}
\rowcolor{CornellReferenceNavy}\color{white}\textbf{Cue / recall prompt} & \color{white}\textbf{Notes}\\\endfirsthead
\rowcolor{CornellReferenceNavy}\color{white}\textbf{Cue / recall prompt} & \color{white}\textbf{Notes (continued)}\\\endhead
\arrayrulecolor{CornellReferenceLineGray}
\textbf{What rule applies to dynamic type?}\\[-1mm]{\scriptsize\CornellClauseRef{3.19} [defns.\allowbreak{}dynamic.\allowbreak{}type.\allowbreak{}prvalue]} & For a prvalue, dynamic type is the same as its statically determined type.\\\hline
\end{longtable}
\begin{CornellReferenceSummaryBox}dynamic type supplies the core-language semantics portion of this clause. The key review move is to connect its local wording to the clause-wide model, then classify each condition by normative force before relying on it.\end{CornellReferenceSummaryBox}
\Needspace{8\baselineskip}
\subsection*{3.20 expression-equivalent}
\begin{longtable}{@{}>{\RaggedRight\arraybackslash\columncolor{CornellReferenceCueGray}}p{0.28\textwidth}>{\RaggedRight\arraybackslash}p{0.64\textwidth}@{}}
\rowcolor{CornellReferenceNavy}\color{white}\textbf{Cue / recall prompt} & \color{white}\textbf{Notes}\\\endfirsthead
\rowcolor{CornellReferenceNavy}\color{white}\textbf{Cue / recall prompt} & \color{white}\textbf{Notes (continued)}\\\endhead
\arrayrulecolor{CornellReferenceLineGray}
\textbf{What rule applies to expression-equivalent?}\\[-1mm]{\scriptsize\CornellClauseRef{3.20} [defns.\allowbreak{}expression.\allowbreak{}equivalent]} & Expression-equivalent expressions agree in effects, potentially-throwing status, and constant-subexpression status.\\\hline
\end{longtable}
\begin{CornellReferenceSummaryBox}expression-equivalent supplies the core-language semantics portion of this clause. The key review move is to connect its local wording to the clause-wide model, then classify each condition by normative force before relying on it.\end{CornellReferenceSummaryBox}
\Needspace{8\baselineskip}
\subsection*{3.21 finite state machine}
\begin{longtable}{@{}>{\RaggedRight\arraybackslash\columncolor{CornellReferenceCueGray}}p{0.28\textwidth}>{\RaggedRight\arraybackslash}p{0.64\textwidth}@{}}
\rowcolor{CornellReferenceNavy}\color{white}\textbf{Cue / recall prompt} & \color{white}\textbf{Notes}\\\endfirsthead
\rowcolor{CornellReferenceNavy}\color{white}\textbf{Cue / recall prompt} & \color{white}\textbf{Notes (continued)}\\\endhead
\arrayrulecolor{CornellReferenceLineGray}
\textbf{What rule applies to finite state machine?}\\[-1mm]{\scriptsize\CornellClauseRef{3.21} [defns.\allowbreak{}regex.\allowbreak{}finite.\allowbreak{}state.\allowbreak{}machine]} & For regular expressions, a finite-state machine is an unspecified representation used to support efficient matching.\\\hline
\end{longtable}
\begin{CornellReferenceSummaryBox}finite state machine supplies the core-language semantics portion of this clause. The key review move is to connect its local wording to the clause-wide model, then classify each condition by normative force before relying on it.\end{CornellReferenceSummaryBox}
\Needspace{8\baselineskip}
\subsection*{3.22 format specifier}
\begin{longtable}{@{}>{\RaggedRight\arraybackslash\columncolor{CornellReferenceCueGray}}p{0.28\textwidth}>{\RaggedRight\arraybackslash}p{0.64\textwidth}@{}}
\rowcolor{CornellReferenceNavy}\color{white}\textbf{Cue / recall prompt} & \color{white}\textbf{Notes}\\\endfirsthead
\rowcolor{CornellReferenceNavy}\color{white}\textbf{Cue / recall prompt} & \color{white}\textbf{Notes (continued)}\\\endhead
\arrayrulecolor{CornellReferenceLineGray}
\textbf{What rule applies to format specifier?}\\[-1mm]{\scriptsize\CornellClauseRef{3.22} [defns.\allowbreak{}regex.\allowbreak{}format.\allowbreak{}specifier]} & A regular-expression format specifier is a character sequence replaced by material derived from a match.\\\hline
\end{longtable}
\begin{CornellReferenceSummaryBox}format specifier supplies the core-language semantics portion of this clause. The key review move is to connect its local wording to the clause-wide model, then classify each condition by normative force before relying on it.\end{CornellReferenceSummaryBox}
\Needspace{8\baselineskip}
\subsection*{3.23 handler function}
\begin{longtable}{@{}>{\RaggedRight\arraybackslash\columncolor{CornellReferenceCueGray}}p{0.28\textwidth}>{\RaggedRight\arraybackslash}p{0.64\textwidth}@{}}
\rowcolor{CornellReferenceNavy}\color{white}\textbf{Cue / recall prompt} & \color{white}\textbf{Notes}\\\endfirsthead
\rowcolor{CornellReferenceNavy}\color{white}\textbf{Cue / recall prompt} & \color{white}\textbf{Notes (continued)}\\\endhead
\arrayrulecolor{CornellReferenceLineGray}
\textbf{What rule applies to handler function?}\\[-1mm]{\scriptsize\CornellClauseRef{3.23} [defns.\allowbreak{}handler]} & A handler function is a non-reserved function supplied by the program and installed through a designated library interface.\\\hline
\end{longtable}
\begin{CornellReferenceSummaryBox}handler function supplies the operation semantics portion of this clause. The key review move is to connect its local wording to the clause-wide model, then classify each condition by normative force before relying on it.\end{CornellReferenceSummaryBox}
\Needspace{8\baselineskip}
\subsection*{3.24 ill-formed program}
\begin{longtable}{@{}>{\RaggedRight\arraybackslash\columncolor{CornellReferenceCueGray}}p{0.28\textwidth}>{\RaggedRight\arraybackslash}p{0.64\textwidth}@{}}
\rowcolor{CornellReferenceNavy}\color{white}\textbf{Cue / recall prompt} & \color{white}\textbf{Notes}\\\endfirsthead
\rowcolor{CornellReferenceNavy}\color{white}\textbf{Cue / recall prompt} & \color{white}\textbf{Notes (continued)}\\\endhead
\arrayrulecolor{CornellReferenceLineGray}
\textbf{What rule applies to ill-formed program?}\\[-1mm]{\scriptsize\CornellClauseRef{3.24} [defns.\allowbreak{}ill.\allowbreak{}formed]} & An ill-formed program fails at least one applicable syntax or semantic rule.\\\hline
\end{longtable}
\begin{CornellReferenceSummaryBox}ill-formed program supplies the core-language semantics portion of this clause. The key review move is to connect its local wording to the clause-wide model, then classify each condition by normative force before relying on it.\end{CornellReferenceSummaryBox}
\Needspace{8\baselineskip}
\subsection*{3.25 implementation-defined behavior}
\begin{longtable}{@{}>{\RaggedRight\arraybackslash\columncolor{CornellReferenceCueGray}}p{0.28\textwidth}>{\RaggedRight\arraybackslash}p{0.64\textwidth}@{}}
\rowcolor{CornellReferenceNavy}\color{white}\textbf{Cue / recall prompt} & \color{white}\textbf{Notes}\\\endfirsthead
\rowcolor{CornellReferenceNavy}\color{white}\textbf{Cue / recall prompt} & \color{white}\textbf{Notes (continued)}\\\endhead
\arrayrulecolor{CornellReferenceLineGray}
\textbf{What rule applies to implementation-defined behavior?}\\[-1mm]{\scriptsize\CornellClauseRef{3.25} [defns.\allowbreak{}impl.\allowbreak{}defined]} & Implementation-defined behavior is a documented implementation choice for a well-formed construct and correct data.\\\hline
\end{longtable}
\begin{CornellReferenceSummaryBox}implementation-defined behavior supplies the core-language semantics portion of this clause. The key review move is to connect its local wording to the clause-wide model, then classify each condition by normative force before relying on it.\end{CornellReferenceSummaryBox}
\Needspace{8\baselineskip}
\subsection*{3.26 implementation-defined strict total order over pointers}
\begin{longtable}{@{}>{\RaggedRight\arraybackslash\columncolor{CornellReferenceCueGray}}p{0.28\textwidth}>{\RaggedRight\arraybackslash}p{0.64\textwidth}@{}}
\rowcolor{CornellReferenceNavy}\color{white}\textbf{Cue / recall prompt} & \color{white}\textbf{Notes}\\\endfirsthead
\rowcolor{CornellReferenceNavy}\color{white}\textbf{Cue / recall prompt} & \color{white}\textbf{Notes (continued)}\\\endhead
\arrayrulecolor{CornellReferenceLineGray}
\textbf{What rule applies to implementation-defined strict total order over pointers?}\\[-1mm]{\scriptsize\CornellClauseRef{3.26} [defns.\allowbreak{}order.\allowbreak{}ptr]} & The library's implementation-defined total order over pointers extends the ordering relations available from built-in pointer comparisons.\\\hline
\end{longtable}
\begin{CornellReferenceSummaryBox}implementation-defined strict total order over pointers supplies the core-language semantics portion of this clause. The key review move is to connect its local wording to the clause-wide model, then classify each condition by normative force before relying on it.\end{CornellReferenceSummaryBox}
\Needspace{8\baselineskip}
\subsection*{3.27 implementation limit}
\begin{longtable}{@{}>{\RaggedRight\arraybackslash\columncolor{CornellReferenceCueGray}}p{0.28\textwidth}>{\RaggedRight\arraybackslash}p{0.64\textwidth}@{}}
\rowcolor{CornellReferenceNavy}\color{white}\textbf{Cue / recall prompt} & \color{white}\textbf{Notes}\\\endfirsthead
\rowcolor{CornellReferenceNavy}\color{white}\textbf{Cue / recall prompt} & \color{white}\textbf{Notes (continued)}\\\endhead
\arrayrulecolor{CornellReferenceLineGray}
\textbf{What rule applies to implementation limit?}\\[-1mm]{\scriptsize\CornellClauseRef{3.27} [defns.\allowbreak{}impl.\allowbreak{}limits]} & An implementation limit is a documented or otherwise specified restriction imposed on programs by an implementation.\\\hline
\end{longtable}
\begin{CornellReferenceSummaryBox}implementation limit supplies the core-language semantics portion of this clause. The key review move is to connect its local wording to the clause-wide model, then classify each condition by normative force before relying on it.\end{CornellReferenceSummaryBox}
\Needspace{8\baselineskip}
\subsection*{3.28 locale-specific behavior}
\begin{longtable}{@{}>{\RaggedRight\arraybackslash\columncolor{CornellReferenceCueGray}}p{0.28\textwidth}>{\RaggedRight\arraybackslash}p{0.64\textwidth}@{}}
\rowcolor{CornellReferenceNavy}\color{white}\textbf{Cue / recall prompt} & \color{white}\textbf{Notes}\\\endfirsthead
\rowcolor{CornellReferenceNavy}\color{white}\textbf{Cue / recall prompt} & \color{white}\textbf{Notes (continued)}\\\endhead
\arrayrulecolor{CornellReferenceLineGray}
\textbf{What rule applies to locale-specific behavior?}\\[-1mm]{\scriptsize\CornellClauseRef{3.28} [defns.\allowbreak{}locale.\allowbreak{}specific]} & Locale-specific behavior is a documented choice controlled by national, cultural, or language conventions.\\\hline
\end{longtable}
\begin{CornellReferenceSummaryBox}locale-specific behavior supplies the core-language semantics portion of this clause. The key review move is to connect its local wording to the clause-wide model, then classify each condition by normative force before relying on it.\end{CornellReferenceSummaryBox}
\Needspace{8\baselineskip}
\subsection*{3.29 matched}
\begin{longtable}{@{}>{\RaggedRight\arraybackslash\columncolor{CornellReferenceCueGray}}p{0.28\textwidth}>{\RaggedRight\arraybackslash}p{0.64\textwidth}@{}}
\rowcolor{CornellReferenceNavy}\color{white}\textbf{Cue / recall prompt} & \color{white}\textbf{Notes}\\\endfirsthead
\rowcolor{CornellReferenceNavy}\color{white}\textbf{Cue / recall prompt} & \color{white}\textbf{Notes (continued)}\\\endhead
\arrayrulecolor{CornellReferenceLineGray}
\textbf{What rule applies to matched?}\\[-1mm]{\scriptsize\CornellClauseRef{3.29} [defns.\allowbreak{}regex.\allowbreak{}matched]} & A sequence is matched when its characters correspond to the selected regular-expression pattern.\\\hline
\end{longtable}
\begin{CornellReferenceSummaryBox}matched supplies the core-language semantics portion of this clause. The key review move is to connect its local wording to the clause-wide model, then classify each condition by normative force before relying on it.\end{CornellReferenceSummaryBox}
\Needspace{8\baselineskip}
\subsection*{3.30 modifier function}
\begin{longtable}{@{}>{\RaggedRight\arraybackslash\columncolor{CornellReferenceCueGray}}p{0.28\textwidth}>{\RaggedRight\arraybackslash}p{0.64\textwidth}@{}}
\rowcolor{CornellReferenceNavy}\color{white}\textbf{Cue / recall prompt} & \color{white}\textbf{Notes}\\\endfirsthead
\rowcolor{CornellReferenceNavy}\color{white}\textbf{Cue / recall prompt} & \color{white}\textbf{Notes (continued)}\\\endhead
\arrayrulecolor{CornellReferenceLineGray}
\textbf{What rule applies to modifier function?}\\[-1mm]{\scriptsize\CornellClauseRef{3.30} [defns.\allowbreak{}modifier]} & A modifier is a non-constructor, non-assignment, non-destructor member function that changes an object's state.\\\hline
\end{longtable}
\begin{CornellReferenceSummaryBox}modifier function supplies the operation semantics portion of this clause. The key review move is to connect its local wording to the clause-wide model, then classify each condition by normative force before relying on it.\end{CornellReferenceSummaryBox}
\Needspace{8\baselineskip}
\subsection*{3.31 move assignment}
\begin{longtable}{@{}>{\RaggedRight\arraybackslash\columncolor{CornellReferenceCueGray}}p{0.28\textwidth}>{\RaggedRight\arraybackslash}p{0.64\textwidth}@{}}
\rowcolor{CornellReferenceNavy}\color{white}\textbf{Cue / recall prompt} & \color{white}\textbf{Notes}\\\endfirsthead
\rowcolor{CornellReferenceNavy}\color{white}\textbf{Cue / recall prompt} & \color{white}\textbf{Notes (continued)}\\\endhead
\arrayrulecolor{CornellReferenceLineGray}
\textbf{What rule applies to move assignment?}\\[-1mm]{\scriptsize\CornellClauseRef{3.31} [defns.\allowbreak{}move.\allowbreak{}assign]} & Move assignment assigns an rvalue object to a modifiable lvalue of the same type under the library's move model.\\\hline
\end{longtable}
\begin{CornellReferenceSummaryBox}move assignment supplies the object contract portion of this clause. The key review move is to connect its local wording to the clause-wide model, then classify each condition by normative force before relying on it.\end{CornellReferenceSummaryBox}
\Needspace{8\baselineskip}
\subsection*{3.32 move construction}
\begin{longtable}{@{}>{\RaggedRight\arraybackslash\columncolor{CornellReferenceCueGray}}p{0.28\textwidth}>{\RaggedRight\arraybackslash}p{0.64\textwidth}@{}}
\rowcolor{CornellReferenceNavy}\color{white}\textbf{Cue / recall prompt} & \color{white}\textbf{Notes}\\\endfirsthead
\rowcolor{CornellReferenceNavy}\color{white}\textbf{Cue / recall prompt} & \color{white}\textbf{Notes (continued)}\\\endhead
\arrayrulecolor{CornellReferenceLineGray}
\textbf{What rule applies to move construction?}\\[-1mm]{\scriptsize\CornellClauseRef{3.32} [defns.\allowbreak{}move.\allowbreak{}constr]} & Move construction directly initializes an object from an rvalue of the same type.\\\hline
\end{longtable}
\begin{CornellReferenceSummaryBox}move construction supplies the core-language semantics portion of this clause. The key review move is to connect its local wording to the clause-wide model, then classify each condition by normative force before relying on it.\end{CornellReferenceSummaryBox}
\Needspace{8\baselineskip}
\subsection*{3.33 non-constant library call}
\begin{longtable}{@{}>{\RaggedRight\arraybackslash\columncolor{CornellReferenceCueGray}}p{0.28\textwidth}>{\RaggedRight\arraybackslash}p{0.64\textwidth}@{}}
\rowcolor{CornellReferenceNavy}\color{white}\textbf{Cue / recall prompt} & \color{white}\textbf{Notes}\\\endfirsthead
\rowcolor{CornellReferenceNavy}\color{white}\textbf{Cue / recall prompt} & \color{white}\textbf{Notes (continued)}\\\endhead
\arrayrulecolor{CornellReferenceLineGray}
\textbf{What rule applies to non-constant library call?}\\[-1mm]{\scriptsize\CornellClauseRef{3.33} [defns.\allowbreak{}nonconst.\allowbreak{}libcall]} & A non-constant library call is an invocation that prevents the containing evaluation from satisfying core constant-expression rules.\\\hline
\end{longtable}
\begin{CornellReferenceSummaryBox}non-constant library call supplies the core-language semantics portion of this clause. The key review move is to connect its local wording to the clause-wide model, then classify each condition by normative force before relying on it.\end{CornellReferenceSummaryBox}
\Needspace{8\baselineskip}
\subsection*{3.34 NTCTS}
\begin{longtable}{@{}>{\RaggedRight\arraybackslash\columncolor{CornellReferenceCueGray}}p{0.28\textwidth}>{\RaggedRight\arraybackslash}p{0.64\textwidth}@{}}
\rowcolor{CornellReferenceNavy}\color{white}\textbf{Cue / recall prompt} & \color{white}\textbf{Notes}\\\endfirsthead
\rowcolor{CornellReferenceNavy}\color{white}\textbf{Cue / recall prompt} & \color{white}\textbf{Notes (continued)}\\\endhead
\arrayrulecolor{CornellReferenceLineGray}
\textbf{What rule applies to NTCTS?}\\[-1mm]{\scriptsize\CornellClauseRef{3.34} [defns.\allowbreak{}ntcts]} & An NTCTS is a character-type sequence ending immediately before the terminating null value of that character type.\\\hline
\end{longtable}
\begin{CornellReferenceSummaryBox}NTCTS supplies the core-language semantics portion of this clause. The key review move is to connect its local wording to the clause-wide model, then classify each condition by normative force before relying on it.\end{CornellReferenceSummaryBox}
\Needspace{8\baselineskip}
\subsection*{3.35 observer function}
\begin{longtable}{@{}>{\RaggedRight\arraybackslash\columncolor{CornellReferenceCueGray}}p{0.28\textwidth}>{\RaggedRight\arraybackslash}p{0.64\textwidth}@{}}
\rowcolor{CornellReferenceNavy}\color{white}\textbf{Cue / recall prompt} & \color{white}\textbf{Notes}\\\endfirsthead
\rowcolor{CornellReferenceNavy}\color{white}\textbf{Cue / recall prompt} & \color{white}\textbf{Notes (continued)}\\\endhead
\arrayrulecolor{CornellReferenceLineGray}
\textbf{What rule applies to observer function?}\\[-1mm]{\scriptsize\CornellClauseRef{3.35} [defns.\allowbreak{}observer]} & An observer reads object state without changing that state and is specified as a const member function.\\\hline
\end{longtable}
\begin{CornellReferenceSummaryBox}observer function supplies the operation semantics portion of this clause. The key review move is to connect its local wording to the clause-wide model, then classify each condition by normative force before relying on it.\end{CornellReferenceSummaryBox}
\Needspace{8\baselineskip}
\subsection*{3.36 parameter}
\begin{longtable}{@{}>{\RaggedRight\arraybackslash\columncolor{CornellReferenceCueGray}}p{0.28\textwidth}>{\RaggedRight\arraybackslash}p{0.64\textwidth}@{}}
\rowcolor{CornellReferenceNavy}\color{white}\textbf{Cue / recall prompt} & \color{white}\textbf{Notes}\\\endfirsthead
\rowcolor{CornellReferenceNavy}\color{white}\textbf{Cue / recall prompt} & \color{white}\textbf{Notes (continued)}\\\endhead
\arrayrulecolor{CornellReferenceLineGray}
\textbf{What rule applies to parameter?}\\[-1mm]{\scriptsize\CornellClauseRef{3.36} [defns.\allowbreak{}parameter]} & A function or handler parameter is the declared object or reference that receives a value on entry.\\\hline
\end{longtable}
\begin{CornellReferenceSummaryBox}parameter supplies the core-language semantics portion of this clause. The key review move is to connect its local wording to the clause-wide model, then classify each condition by normative force before relying on it.\end{CornellReferenceSummaryBox}
\Needspace{8\baselineskip}
\subsection*{3.37 parameter}
\begin{longtable}{@{}>{\RaggedRight\arraybackslash\columncolor{CornellReferenceCueGray}}p{0.28\textwidth}>{\RaggedRight\arraybackslash}p{0.64\textwidth}@{}}
\rowcolor{CornellReferenceNavy}\color{white}\textbf{Cue / recall prompt} & \color{white}\textbf{Notes}\\\endfirsthead
\rowcolor{CornellReferenceNavy}\color{white}\textbf{Cue / recall prompt} & \color{white}\textbf{Notes (continued)}\\\endhead
\arrayrulecolor{CornellReferenceLineGray}
\textbf{What rule applies to parameter?}\\[-1mm]{\scriptsize\CornellClauseRef{3.37} [defns.\allowbreak{}parameter.\allowbreak{}macro]} & A macro parameter is an identifier in the parameter list of a function-like macro definition.\\\hline
\end{longtable}
\begin{CornellReferenceSummaryBox}parameter supplies the core-language semantics portion of this clause. The key review move is to connect its local wording to the clause-wide model, then classify each condition by normative force before relying on it.\end{CornellReferenceSummaryBox}
\Needspace{8\baselineskip}
\subsection*{3.38 parameter}
\begin{longtable}{@{}>{\RaggedRight\arraybackslash\columncolor{CornellReferenceCueGray}}p{0.28\textwidth}>{\RaggedRight\arraybackslash}p{0.64\textwidth}@{}}
\rowcolor{CornellReferenceNavy}\color{white}\textbf{Cue / recall prompt} & \color{white}\textbf{Notes}\\\endfirsthead
\rowcolor{CornellReferenceNavy}\color{white}\textbf{Cue / recall prompt} & \color{white}\textbf{Notes (continued)}\\\endhead
\arrayrulecolor{CornellReferenceLineGray}
\textbf{What rule applies to parameter?}\\[-1mm]{\scriptsize\CornellClauseRef{3.38} [defns.\allowbreak{}parameter.\allowbreak{}templ]} & A template parameter is one member of a template-parameter-list.\\\hline
\end{longtable}
\begin{CornellReferenceSummaryBox}parameter supplies the core-language semantics portion of this clause. The key review move is to connect its local wording to the clause-wide model, then classify each condition by normative force before relying on it.\end{CornellReferenceSummaryBox}
\Needspace{8\baselineskip}
\subsection*{3.39 primary equivalence class}
\begin{longtable}{@{}>{\RaggedRight\arraybackslash\columncolor{CornellReferenceCueGray}}p{0.28\textwidth}>{\RaggedRight\arraybackslash}p{0.64\textwidth}@{}}
\rowcolor{CornellReferenceNavy}\color{white}\textbf{Cue / recall prompt} & \color{white}\textbf{Notes}\\\endfirsthead
\rowcolor{CornellReferenceNavy}\color{white}\textbf{Cue / recall prompt} & \color{white}\textbf{Notes (continued)}\\\endhead
\arrayrulecolor{CornellReferenceLineGray}
\textbf{What rule applies to primary equivalence class?}\\[-1mm]{\scriptsize\CornellClauseRef{3.39} [defns.\allowbreak{}regex.\allowbreak{}primary.\allowbreak{}equivalence.\allowbreak{}class]} & A primary equivalence class groups characters sharing the primary collation key before distinctions such as accents or case.\\\hline
\end{longtable}
\begin{CornellReferenceSummaryBox}primary equivalence class supplies the core-language semantics portion of this clause. The key review move is to connect its local wording to the clause-wide model, then classify each condition by normative force before relying on it.\end{CornellReferenceSummaryBox}
\Needspace{8\baselineskip}
\subsection*{3.40 program-defined specialization}
\begin{longtable}{@{}>{\RaggedRight\arraybackslash\columncolor{CornellReferenceCueGray}}p{0.28\textwidth}>{\RaggedRight\arraybackslash}p{0.64\textwidth}@{}}
\rowcolor{CornellReferenceNavy}\color{white}\textbf{Cue / recall prompt} & \color{white}\textbf{Notes}\\\endfirsthead
\rowcolor{CornellReferenceNavy}\color{white}\textbf{Cue / recall prompt} & \color{white}\textbf{Notes (continued)}\\\endhead
\arrayrulecolor{CornellReferenceLineGray}
\textbf{What rule applies to program-defined specialization?}\\[-1mm]{\scriptsize\CornellClauseRef{3.40} [defns.\allowbreak{}prog.\allowbreak{}def.\allowbreak{}spec]} & A program-defined specialization is an explicit or partial specialization supplied by the program rather than the library or implementation.\\\hline
\end{longtable}
\begin{CornellReferenceSummaryBox}program-defined specialization supplies the core-language semantics portion of this clause. The key review move is to connect its local wording to the clause-wide model, then classify each condition by normative force before relying on it.\end{CornellReferenceSummaryBox}
\Needspace{8\baselineskip}
\subsection*{3.41 program-defined type}
\begin{longtable}{@{}>{\RaggedRight\arraybackslash\columncolor{CornellReferenceCueGray}}p{0.28\textwidth}>{\RaggedRight\arraybackslash}p{0.64\textwidth}@{}}
\rowcolor{CornellReferenceNavy}\color{white}\textbf{Cue / recall prompt} & \color{white}\textbf{Notes}\\\endfirsthead
\rowcolor{CornellReferenceNavy}\color{white}\textbf{Cue / recall prompt} & \color{white}\textbf{Notes (continued)}\\\endhead
\arrayrulecolor{CornellReferenceLineGray}
\textbf{What rule applies to program-defined type?}\\[-1mm]{\scriptsize\CornellClauseRef{3.41} [defns.\allowbreak{}prog.\allowbreak{}def.\allowbreak{}type]} & A program-defined type is a qualifying class, enumeration, closure, or specialization type not supplied by the library or implementation.\\\hline
\end{longtable}
\begin{CornellReferenceSummaryBox}program-defined type supplies the core-language semantics portion of this clause. The key review move is to connect its local wording to the clause-wide model, then classify each condition by normative force before relying on it.\end{CornellReferenceSummaryBox}
\Needspace{8\baselineskip}
\subsection*{3.42 projection}
\begin{longtable}{@{}>{\RaggedRight\arraybackslash\columncolor{CornellReferenceCueGray}}p{0.28\textwidth}>{\RaggedRight\arraybackslash}p{0.64\textwidth}@{}}
\rowcolor{CornellReferenceNavy}\color{white}\textbf{Cue / recall prompt} & \color{white}\textbf{Notes}\\\endfirsthead
\rowcolor{CornellReferenceNavy}\color{white}\textbf{Cue / recall prompt} & \color{white}\textbf{Notes (continued)}\\\endhead
\arrayrulecolor{CornellReferenceLineGray}
\textbf{What rule applies to projection?}\\[-1mm]{\scriptsize\CornellClauseRef{3.42} [defns.\allowbreak{}projection]} & A projection is the transformation an algorithm applies before inspecting or comparing an element.\\\hline
\end{longtable}
\begin{CornellReferenceSummaryBox}projection supplies the core-language semantics portion of this clause. The key review move is to connect its local wording to the clause-wide model, then classify each condition by normative force before relying on it.\end{CornellReferenceSummaryBox}
\Needspace{8\baselineskip}
\subsection*{3.43 referenceable type}
\begin{longtable}{@{}>{\RaggedRight\arraybackslash\columncolor{CornellReferenceCueGray}}p{0.28\textwidth}>{\RaggedRight\arraybackslash}p{0.64\textwidth}@{}}
\rowcolor{CornellReferenceNavy}\color{white}\textbf{Cue / recall prompt} & \color{white}\textbf{Notes}\\\endfirsthead
\rowcolor{CornellReferenceNavy}\color{white}\textbf{Cue / recall prompt} & \color{white}\textbf{Notes (continued)}\\\endhead
\arrayrulecolor{CornellReferenceLineGray}
\textbf{What rule applies to referenceable type?}\\[-1mm]{\scriptsize\CornellClauseRef{3.43} [defns.\allowbreak{}referenceable]} & A referenceable type is an object type, an eligible function type, or a reference type for which a reference can be formed under the rules.\\\hline
\end{longtable}
\begin{CornellReferenceSummaryBox}referenceable type supplies the core-language semantics portion of this clause. The key review move is to connect its local wording to the clause-wide model, then classify each condition by normative force before relying on it.\end{CornellReferenceSummaryBox}
\Needspace{8\baselineskip}
\subsection*{3.44 regular expression}
\begin{longtable}{@{}>{\RaggedRight\arraybackslash\columncolor{CornellReferenceCueGray}}p{0.28\textwidth}>{\RaggedRight\arraybackslash}p{0.64\textwidth}@{}}
\rowcolor{CornellReferenceNavy}\color{white}\textbf{Cue / recall prompt} & \color{white}\textbf{Notes}\\\endfirsthead
\rowcolor{CornellReferenceNavy}\color{white}\textbf{Cue / recall prompt} & \color{white}\textbf{Notes (continued)}\\\endhead
\arrayrulecolor{CornellReferenceLineGray}
\textbf{What rule applies to regular expression?}\\[-1mm]{\scriptsize\CornellClauseRef{3.44} [defns.\allowbreak{}regex.\allowbreak{}regular.\allowbreak{}expression]} & A regular expression is a pattern selecting character strings from a set of candidate strings.\\\hline
\end{longtable}
\begin{CornellReferenceSummaryBox}regular expression supplies the core-language semantics portion of this clause. The key review move is to connect its local wording to the clause-wide model, then classify each condition by normative force before relying on it.\end{CornellReferenceSummaryBox}
\Needspace{8\baselineskip}
\subsection*{3.45 replacement function}
\begin{longtable}{@{}>{\RaggedRight\arraybackslash\columncolor{CornellReferenceCueGray}}p{0.28\textwidth}>{\RaggedRight\arraybackslash}p{0.64\textwidth}@{}}
\rowcolor{CornellReferenceNavy}\color{white}\textbf{Cue / recall prompt} & \color{white}\textbf{Notes}\\\endfirsthead
\rowcolor{CornellReferenceNavy}\color{white}\textbf{Cue / recall prompt} & \color{white}\textbf{Notes (continued)}\\\endhead
\arrayrulecolor{CornellReferenceLineGray}
\textbf{What rule applies to replacement function?}\\[-1mm]{\scriptsize\CornellClauseRef{3.45} [defns.\allowbreak{}replacement]} & A replacement function is a non-reserved function definition supplied by the program in place of the implementation's default facility.\\\hline
\end{longtable}
\begin{CornellReferenceSummaryBox}replacement function supplies the operation semantics portion of this clause. The key review move is to connect its local wording to the clause-wide model, then classify each condition by normative force before relying on it.\end{CornellReferenceSummaryBox}
\Needspace{8\baselineskip}
\subsection*{3.46 required behavior}
\begin{longtable}{@{}>{\RaggedRight\arraybackslash\columncolor{CornellReferenceCueGray}}p{0.28\textwidth}>{\RaggedRight\arraybackslash}p{0.64\textwidth}@{}}
\rowcolor{CornellReferenceNavy}\color{white}\textbf{Cue / recall prompt} & \color{white}\textbf{Notes}\\\endfirsthead
\rowcolor{CornellReferenceNavy}\color{white}\textbf{Cue / recall prompt} & \color{white}\textbf{Notes (continued)}\\\endhead
\arrayrulecolor{CornellReferenceLineGray}
\textbf{What rule applies to required behavior?}\\[-1mm]{\scriptsize\CornellClauseRef{3.46} [defns.\allowbreak{}required.\allowbreak{}behavior]} & Required behavior is the common semantic contract that both the implementation default and a program-supplied replacement or handler must satisfy.\\\hline
\end{longtable}
\begin{CornellReferenceSummaryBox}required behavior supplies the core-language semantics portion of this clause. The key review move is to connect its local wording to the clause-wide model, then classify each condition by normative force before relying on it.\end{CornellReferenceSummaryBox}
\Needspace{8\baselineskip}
\subsection*{3.47 reserved function}
\begin{longtable}{@{}>{\RaggedRight\arraybackslash\columncolor{CornellReferenceCueGray}}p{0.28\textwidth}>{\RaggedRight\arraybackslash}p{0.64\textwidth}@{}}
\rowcolor{CornellReferenceNavy}\color{white}\textbf{Cue / recall prompt} & \color{white}\textbf{Notes}\\\endfirsthead
\rowcolor{CornellReferenceNavy}\color{white}\textbf{Cue / recall prompt} & \color{white}\textbf{Notes (continued)}\\\endhead
\arrayrulecolor{CornellReferenceLineGray}
\textbf{What rule applies to reserved function?}\\[-1mm]{\scriptsize\CornellClauseRef{3.47} [defns.\allowbreak{}reserved.\allowbreak{}function]} & A reserved function is a library function defined by the implementation; a program definition of one has the stated undefined consequence.\\\hline
\end{longtable}
\begin{CornellReferenceSummaryBox}reserved function supplies the operation semantics portion of this clause. The key review move is to connect its local wording to the clause-wide model, then classify each condition by normative force before relying on it.\end{CornellReferenceSummaryBox}
\Needspace{8\baselineskip}
\subsection*{3.48 signature}
\begin{longtable}{@{}>{\RaggedRight\arraybackslash\columncolor{CornellReferenceCueGray}}p{0.28\textwidth}>{\RaggedRight\arraybackslash}p{0.64\textwidth}@{}}
\rowcolor{CornellReferenceNavy}\color{white}\textbf{Cue / recall prompt} & \color{white}\textbf{Notes}\\\endfirsthead
\rowcolor{CornellReferenceNavy}\color{white}\textbf{Cue / recall prompt} & \color{white}\textbf{Notes (continued)}\\\endhead
\arrayrulecolor{CornellReferenceLineGray}
\textbf{What rule applies to signature?}\\[-1mm]{\scriptsize\CornellClauseRef{3.48} [defns.\allowbreak{}signature]} & A non-member function signature combines its name, parameter-type-list, and enclosing namespace.\\\hline
\end{longtable}
\begin{CornellReferenceSummaryBox}signature supplies the core-language semantics portion of this clause. The key review move is to connect its local wording to the clause-wide model, then classify each condition by normative force before relying on it.\end{CornellReferenceSummaryBox}
\Needspace{8\baselineskip}
\subsection*{3.49 signature}
\begin{longtable}{@{}>{\RaggedRight\arraybackslash\columncolor{CornellReferenceCueGray}}p{0.28\textwidth}>{\RaggedRight\arraybackslash}p{0.64\textwidth}@{}}
\rowcolor{CornellReferenceNavy}\color{white}\textbf{Cue / recall prompt} & \color{white}\textbf{Notes}\\\endfirsthead
\rowcolor{CornellReferenceNavy}\color{white}\textbf{Cue / recall prompt} & \color{white}\textbf{Notes (continued)}\\\endhead
\arrayrulecolor{CornellReferenceLineGray}
\textbf{What rule applies to signature?}\\[-1mm]{\scriptsize\CornellClauseRef{3.49} [defns.\allowbreak{}signature.\allowbreak{}friend]} & A constrained non-template friend signature also includes the enclosing class and trailing constraint.\\\hline
\end{longtable}
\begin{CornellReferenceSummaryBox}signature supplies the core-language semantics portion of this clause. The key review move is to connect its local wording to the clause-wide model, then classify each condition by normative force before relying on it.\end{CornellReferenceSummaryBox}
\Needspace{8\baselineskip}
\subsection*{3.50 signature}
\begin{longtable}{@{}>{\RaggedRight\arraybackslash\columncolor{CornellReferenceCueGray}}p{0.28\textwidth}>{\RaggedRight\arraybackslash}p{0.64\textwidth}@{}}
\rowcolor{CornellReferenceNavy}\color{white}\textbf{Cue / recall prompt} & \color{white}\textbf{Notes}\\\endfirsthead
\rowcolor{CornellReferenceNavy}\color{white}\textbf{Cue / recall prompt} & \color{white}\textbf{Notes (continued)}\\\endhead
\arrayrulecolor{CornellReferenceLineGray}
\textbf{What rule applies to signature?}\\[-1mm]{\scriptsize\CornellClauseRef{3.50} [defns.\allowbreak{}signature.\allowbreak{}templ]} & A function-template signature includes name, parameter types, namespace, return type, template-head signature, and any trailing constraint.\\\hline
\end{longtable}
\begin{CornellReferenceSummaryBox}signature supplies the core-language semantics portion of this clause. The key review move is to connect its local wording to the clause-wide model, then classify each condition by normative force before relying on it.\end{CornellReferenceSummaryBox}
\Needspace{8\baselineskip}
\subsection*{3.51 signature}
\begin{longtable}{@{}>{\RaggedRight\arraybackslash\columncolor{CornellReferenceCueGray}}p{0.28\textwidth}>{\RaggedRight\arraybackslash}p{0.64\textwidth}@{}}
\rowcolor{CornellReferenceNavy}\color{white}\textbf{Cue / recall prompt} & \color{white}\textbf{Notes}\\\endfirsthead
\rowcolor{CornellReferenceNavy}\color{white}\textbf{Cue / recall prompt} & \color{white}\textbf{Notes (continued)}\\\endhead
\arrayrulecolor{CornellReferenceLineGray}
\textbf{What rule applies to signature?}\\[-1mm]{\scriptsize\CornellClauseRef{3.51} [defns.\allowbreak{}signature.\allowbreak{}templ.\allowbreak{}friend]} & A constrained friend function-template signature additionally identifies its enclosing class context.\\\hline
\end{longtable}
\begin{CornellReferenceSummaryBox}signature supplies the core-language semantics portion of this clause. The key review move is to connect its local wording to the clause-wide model, then classify each condition by normative force before relying on it.\end{CornellReferenceSummaryBox}
\Needspace{8\baselineskip}
\subsection*{3.52 signature}
\begin{longtable}{@{}>{\RaggedRight\arraybackslash\columncolor{CornellReferenceCueGray}}p{0.28\textwidth}>{\RaggedRight\arraybackslash}p{0.64\textwidth}@{}}
\rowcolor{CornellReferenceNavy}\color{white}\textbf{Cue / recall prompt} & \color{white}\textbf{Notes}\\\endfirsthead
\rowcolor{CornellReferenceNavy}\color{white}\textbf{Cue / recall prompt} & \color{white}\textbf{Notes (continued)}\\\endhead
\arrayrulecolor{CornellReferenceLineGray}
\textbf{What rule applies to signature?}\\[-1mm]{\scriptsize\CornellClauseRef{3.52} [defns.\allowbreak{}signature.\allowbreak{}spec]} & A function-template specialization signature combines the primary template's signature with the specialization's explicit or deduced arguments.\\\hline
\end{longtable}
\begin{CornellReferenceSummaryBox}signature supplies the core-language semantics portion of this clause. The key review move is to connect its local wording to the clause-wide model, then classify each condition by normative force before relying on it.\end{CornellReferenceSummaryBox}
\Needspace{8\baselineskip}
\subsection*{3.53 signature}
\begin{longtable}{@{}>{\RaggedRight\arraybackslash\columncolor{CornellReferenceCueGray}}p{0.28\textwidth}>{\RaggedRight\arraybackslash}p{0.64\textwidth}@{}}
\rowcolor{CornellReferenceNavy}\color{white}\textbf{Cue / recall prompt} & \color{white}\textbf{Notes}\\\endfirsthead
\rowcolor{CornellReferenceNavy}\color{white}\textbf{Cue / recall prompt} & \color{white}\textbf{Notes (continued)}\\\endhead
\arrayrulecolor{CornellReferenceLineGray}
\textbf{What rule applies to signature?}\\[-1mm]{\scriptsize\CornellClauseRef{3.53} [defns.\allowbreak{}signature.\allowbreak{}member]} & A member-function signature adds its class, cv-qualification, ref-qualification, and trailing constraint to the name and parameter types.\\\hline
\end{longtable}
\begin{CornellReferenceSummaryBox}signature supplies the core-language semantics portion of this clause. The key review move is to connect its local wording to the clause-wide model, then classify each condition by normative force before relying on it.\end{CornellReferenceSummaryBox}
\Needspace{8\baselineskip}
\subsection*{3.54 signature}
\begin{longtable}{@{}>{\RaggedRight\arraybackslash\columncolor{CornellReferenceCueGray}}p{0.28\textwidth}>{\RaggedRight\arraybackslash}p{0.64\textwidth}@{}}
\rowcolor{CornellReferenceNavy}\color{white}\textbf{Cue / recall prompt} & \color{white}\textbf{Notes}\\\endfirsthead
\rowcolor{CornellReferenceNavy}\color{white}\textbf{Cue / recall prompt} & \color{white}\textbf{Notes (continued)}\\\endhead
\arrayrulecolor{CornellReferenceLineGray}
\textbf{What rule applies to signature?}\\[-1mm]{\scriptsize\CornellClauseRef{3.54} [defns.\allowbreak{}signature.\allowbreak{}member.\allowbreak{}templ]} & A member-function-template signature also includes return type when present and the template-head signature.\\\hline
\end{longtable}
\begin{CornellReferenceSummaryBox}signature supplies the core-language semantics portion of this clause. The key review move is to connect its local wording to the clause-wide model, then classify each condition by normative force before relying on it.\end{CornellReferenceSummaryBox}
\Needspace{8\baselineskip}
\subsection*{3.55 signature}
\begin{longtable}{@{}>{\RaggedRight\arraybackslash\columncolor{CornellReferenceCueGray}}p{0.28\textwidth}>{\RaggedRight\arraybackslash}p{0.64\textwidth}@{}}
\rowcolor{CornellReferenceNavy}\color{white}\textbf{Cue / recall prompt} & \color{white}\textbf{Notes}\\\endfirsthead
\rowcolor{CornellReferenceNavy}\color{white}\textbf{Cue / recall prompt} & \color{white}\textbf{Notes (continued)}\\\endhead
\arrayrulecolor{CornellReferenceLineGray}
\textbf{What rule applies to signature?}\\[-1mm]{\scriptsize\CornellClauseRef{3.55} [defns.\allowbreak{}signature.\allowbreak{}member.\allowbreak{}spec]} & A member-function-template specialization signature combines the member template's signature with its template arguments.\\\hline
\end{longtable}
\begin{CornellReferenceSummaryBox}signature supplies the core-language semantics portion of this clause. The key review move is to connect its local wording to the clause-wide model, then classify each condition by normative force before relying on it.\end{CornellReferenceSummaryBox}
\Needspace{8\baselineskip}
\subsection*{3.56 signature}
\begin{longtable}{@{}>{\RaggedRight\arraybackslash\columncolor{CornellReferenceCueGray}}p{0.28\textwidth}>{\RaggedRight\arraybackslash}p{0.64\textwidth}@{}}
\rowcolor{CornellReferenceNavy}\color{white}\textbf{Cue / recall prompt} & \color{white}\textbf{Notes}\\\endfirsthead
\rowcolor{CornellReferenceNavy}\color{white}\textbf{Cue / recall prompt} & \color{white}\textbf{Notes (continued)}\\\endhead
\arrayrulecolor{CornellReferenceLineGray}
\textbf{What rule applies to signature?}\\[-1mm]{\scriptsize\CornellClauseRef{3.56} [defns.\allowbreak{}signature.\allowbreak{}template.\allowbreak{}head]} & A template-head signature consists of parameter kinds and structure, excluding parameter names and defaults, plus any requires-clause.\\\hline
\end{longtable}
\begin{CornellReferenceSummaryBox}signature supplies the core-language semantics portion of this clause. The key review move is to connect its local wording to the clause-wide model, then classify each condition by normative force before relying on it.\end{CornellReferenceSummaryBox}
\Needspace{8\baselineskip}
\subsection*{3.57 stable algorithm}
\begin{longtable}{@{}>{\RaggedRight\arraybackslash\columncolor{CornellReferenceCueGray}}p{0.28\textwidth}>{\RaggedRight\arraybackslash}p{0.64\textwidth}@{}}
\rowcolor{CornellReferenceNavy}\color{white}\textbf{Cue / recall prompt} & \color{white}\textbf{Notes}\\\endfirsthead
\rowcolor{CornellReferenceNavy}\color{white}\textbf{Cue / recall prompt} & \color{white}\textbf{Notes (continued)}\\\endhead
\arrayrulecolor{CornellReferenceLineGray}
\textbf{What does stable algorithm guarantee?}\\[-1mm]{\scriptsize\CornellClauseRef{3.57} [defns.\allowbreak{}stable]} & A stable algorithm preserves the relative order required by that algorithm's contract.\\\hline
\end{longtable}
\begin{CornellReferenceSummaryBox}stable algorithm supplies the operation semantics portion of this clause. The key review move is to connect its local wording to the clause-wide model, then classify each condition by normative force before relying on it.\end{CornellReferenceSummaryBox}
\Needspace{8\baselineskip}
\subsection*{3.58 static type}
\begin{longtable}{@{}>{\RaggedRight\arraybackslash\columncolor{CornellReferenceCueGray}}p{0.28\textwidth}>{\RaggedRight\arraybackslash}p{0.64\textwidth}@{}}
\rowcolor{CornellReferenceNavy}\color{white}\textbf{Cue / recall prompt} & \color{white}\textbf{Notes}\\\endfirsthead
\rowcolor{CornellReferenceNavy}\color{white}\textbf{Cue / recall prompt} & \color{white}\textbf{Notes (continued)}\\\endhead
\arrayrulecolor{CornellReferenceLineGray}
\textbf{What rule applies to static type?}\\[-1mm]{\scriptsize\CornellClauseRef{3.58} [defns.\allowbreak{}static.\allowbreak{}type]} & Static type is determined from program form without considering runtime execution and does not change during execution.\\\hline
\end{longtable}
\begin{CornellReferenceSummaryBox}static type supplies the core-language semantics portion of this clause. The key review move is to connect its local wording to the clause-wide model, then classify each condition by normative force before relying on it.\end{CornellReferenceSummaryBox}
\Needspace{8\baselineskip}
\subsection*{3.59 sub-expression}
\begin{longtable}{@{}>{\RaggedRight\arraybackslash\columncolor{CornellReferenceCueGray}}p{0.28\textwidth}>{\RaggedRight\arraybackslash}p{0.64\textwidth}@{}}
\rowcolor{CornellReferenceNavy}\color{white}\textbf{Cue / recall prompt} & \color{white}\textbf{Notes}\\\endfirsthead
\rowcolor{CornellReferenceNavy}\color{white}\textbf{Cue / recall prompt} & \color{white}\textbf{Notes (continued)}\\\endhead
\arrayrulecolor{CornellReferenceLineGray}
\textbf{What rule applies to sub-expression?}\\[-1mm]{\scriptsize\CornellClauseRef{3.59} [defns.\allowbreak{}regex.\allowbreak{}subexpression]} & A regular-expression sub-expression is a parenthesized portion of a pattern.\\\hline
\end{longtable}
\begin{CornellReferenceSummaryBox}sub-expression supplies the core-language semantics portion of this clause. The key review move is to connect its local wording to the clause-wide model, then classify each condition by normative force before relying on it.\end{CornellReferenceSummaryBox}
\Needspace{8\baselineskip}
\subsection*{3.60 traits class}
\begin{longtable}{@{}>{\RaggedRight\arraybackslash\columncolor{CornellReferenceCueGray}}p{0.28\textwidth}>{\RaggedRight\arraybackslash}p{0.64\textwidth}@{}}
\rowcolor{CornellReferenceNavy}\color{white}\textbf{Cue / recall prompt} & \color{white}\textbf{Notes}\\\endfirsthead
\rowcolor{CornellReferenceNavy}\color{white}\textbf{Cue / recall prompt} & \color{white}\textbf{Notes (continued)}\\\endhead
\arrayrulecolor{CornellReferenceLineGray}
\textbf{What rule applies to traits class?}\\[-1mm]{\scriptsize\CornellClauseRef{3.60} [defns.\allowbreak{}traits]} & A traits class packages types and operations so templates can adapt to the types with which they are instantiated.\\\hline
\end{longtable}
\begin{CornellReferenceSummaryBox}traits class supplies the core-language semantics portion of this clause. The key review move is to connect its local wording to the clause-wide model, then classify each condition by normative force before relying on it.\end{CornellReferenceSummaryBox}
\Needspace{8\baselineskip}
\subsection*{3.61 unblock}
\begin{longtable}{@{}>{\RaggedRight\arraybackslash\columncolor{CornellReferenceCueGray}}p{0.28\textwidth}>{\RaggedRight\arraybackslash}p{0.64\textwidth}@{}}
\rowcolor{CornellReferenceNavy}\color{white}\textbf{Cue / recall prompt} & \color{white}\textbf{Notes}\\\endfirsthead
\rowcolor{CornellReferenceNavy}\color{white}\textbf{Cue / recall prompt} & \color{white}\textbf{Notes (continued)}\\\endhead
\arrayrulecolor{CornellReferenceLineGray}
\textbf{What rule applies to unblock?}\\[-1mm]{\scriptsize\CornellClauseRef{3.61} [defns.\allowbreak{}unblock]} & To unblock is to satisfy a condition awaited by one or more blocked threads.\\\hline
\end{longtable}
\begin{CornellReferenceSummaryBox}unblock supplies the concurrency contract portion of this clause. The key review move is to connect its local wording to the clause-wide model, then classify each condition by normative force before relying on it.\end{CornellReferenceSummaryBox}
\Needspace{8\baselineskip}
\subsection*{3.62 undefined behavior}
\begin{longtable}{@{}>{\RaggedRight\arraybackslash\columncolor{CornellReferenceCueGray}}p{0.28\textwidth}>{\RaggedRight\arraybackslash}p{0.64\textwidth}@{}}
\rowcolor{CornellReferenceNavy}\color{white}\textbf{Cue / recall prompt} & \color{white}\textbf{Notes}\\\endfirsthead
\rowcolor{CornellReferenceNavy}\color{white}\textbf{Cue / recall prompt} & \color{white}\textbf{Notes (continued)}\\\endhead
\arrayrulecolor{CornellReferenceLineGray}
\textbf{What rule applies to undefined behavior?}\\[-1mm]{\scriptsize\CornellClauseRef{3.62} [defns.\allowbreak{}undefined]} & Undefined behavior imposes no requirements on an implementation; a diagnostic is not generally guaranteed.\\\hline
\end{longtable}
\begin{CornellReferenceSummaryBox}undefined behavior supplies the core-language semantics portion of this clause. The key review move is to connect its local wording to the clause-wide model, then classify each condition by normative force before relying on it.\end{CornellReferenceSummaryBox}
\Needspace{8\baselineskip}
\subsection*{3.63 unspecified behavior}
\begin{longtable}{@{}>{\RaggedRight\arraybackslash\columncolor{CornellReferenceCueGray}}p{0.28\textwidth}>{\RaggedRight\arraybackslash}p{0.64\textwidth}@{}}
\rowcolor{CornellReferenceNavy}\color{white}\textbf{Cue / recall prompt} & \color{white}\textbf{Notes}\\\endfirsthead
\rowcolor{CornellReferenceNavy}\color{white}\textbf{Cue / recall prompt} & \color{white}\textbf{Notes (continued)}\\\endhead
\arrayrulecolor{CornellReferenceLineGray}
\textbf{What rule applies to unspecified behavior?}\\[-1mm]{\scriptsize\CornellClauseRef{3.63} [defns.\allowbreak{}unspecified]} & Unspecified behavior is a bounded implementation choice that need not be documented.\\\hline
\end{longtable}
\begin{CornellReferenceSummaryBox}unspecified behavior supplies the core-language semantics portion of this clause. The key review move is to connect its local wording to the clause-wide model, then classify each condition by normative force before relying on it.\end{CornellReferenceSummaryBox}
\Needspace{8\baselineskip}
\subsection*{3.64 valid but unspecified state}
\begin{longtable}{@{}>{\RaggedRight\arraybackslash\columncolor{CornellReferenceCueGray}}p{0.28\textwidth}>{\RaggedRight\arraybackslash}p{0.64\textwidth}@{}}
\rowcolor{CornellReferenceNavy}\color{white}\textbf{Cue / recall prompt} & \color{white}\textbf{Notes}\\\endfirsthead
\rowcolor{CornellReferenceNavy}\color{white}\textbf{Cue / recall prompt} & \color{white}\textbf{Notes (continued)}\\\endhead
\arrayrulecolor{CornellReferenceLineGray}
\textbf{What rule applies to valid but unspecified state?}\\[-1mm]{\scriptsize\CornellClauseRef{3.64} [defns.\allowbreak{}valid]} & A valid but unspecified state preserves the type's invariants and supports operations whose preconditions can still be established.\\\hline
\end{longtable}
\begin{CornellReferenceSummaryBox}valid but unspecified state supplies the core-language semantics portion of this clause. The key review move is to connect its local wording to the clause-wide model, then classify each condition by normative force before relying on it.\end{CornellReferenceSummaryBox}
\Needspace{8\baselineskip}
\subsection*{3.65 well-formed program}
\begin{longtable}{@{}>{\RaggedRight\arraybackslash\columncolor{CornellReferenceCueGray}}p{0.28\textwidth}>{\RaggedRight\arraybackslash}p{0.64\textwidth}@{}}
\rowcolor{CornellReferenceNavy}\color{white}\textbf{Cue / recall prompt} & \color{white}\textbf{Notes}\\\endfirsthead
\rowcolor{CornellReferenceNavy}\color{white}\textbf{Cue / recall prompt} & \color{white}\textbf{Notes (continued)}\\\endhead
\arrayrulecolor{CornellReferenceLineGray}
\textbf{What rule applies to well-formed program?}\\[-1mm]{\scriptsize\CornellClauseRef{3.65} [defns.\allowbreak{}well.\allowbreak{}formed]} & A well-formed program is constructed in accordance with all applicable syntax and semantic rules.\\\hline
\end{longtable}
\begin{CornellReferenceSummaryBox}well-formed program supplies the core-language semantics portion of this clause. The key review move is to connect its local wording to the clause-wide model, then classify each condition by normative force before relying on it.\end{CornellReferenceSummaryBox}
\section{Language-Lawyer and Library-Usage Pitfalls}
\begin{CornellReferencePitfallBox}\begin{itemize}[label=\textcolor{CornellReferenceAmber}{$\blacktriangleright$}]
\item Using 'undefined', 'unspecified', and 'implementation-defined' interchangeably.
\item Confusing an argument with a parameter.
\item Assuming every library observer is free of indirect effects.
\item Treating a valid but unspecified object as unusable.
\item Treating a note, example, or recommended practice as though it imposed a mandatory program rule.
\end{itemize}\end{CornellReferencePitfallBox}
\section{Programmer and Implementation Checklist}
\begin{CornellReferenceChecklistBox}\begin{itemize}[label=$\square$]
\item Classify each rule using the defined behavior category.
\item Apply context-qualified definitions such as argument, block, dynamic type, and signature precisely.
\item Check whether a library term imposes ownership, state, ordering, or documentation consequences.
\item Follow cross-references when a definition delegates operational details.
\item For \CornellClauseRef{3.1}, verify access against the precise rule category and all cited dependencies.
\item For \CornellClauseRef{3.2}, verify argument against the precise rule category and all cited dependencies.
\item For \CornellClauseRef{3.3}, verify argument against the precise rule category and all cited dependencies.
\item For \CornellClauseRef{3.4}, verify argument against the precise rule category and all cited dependencies.
\end{itemize}\end{CornellReferenceChecklistBox}
\section{Clause Synthesis}
\begin{CornellReferenceOrientationBox}Defines cross-cutting vocabulary used by the core language and library. The definitions distinguish program status and behavior categories, execution and concurrency concepts, library roles, regular-expression terms, and the components of signatures. Provides the shared semantic vocabulary needed to read requirements consistently. A sound review distinguishes syntax and participation from runtime preconditions, keeps notes and examples non-mandatory, and follows cross-clause dependencies before making a portability claim.\end{CornellReferenceOrientationBox}
\section{Self-Test Questions}
\begin{enumerate}
\item What role does access play in the clause's overall model?
\item Which normative distinctions must be kept separate when applying argument?
\item How would you audit a program or implementation that depends on argument?
\item Compare argument with argument; what different question does each answer?
\item What role does argument play in the clause's overall model?
\item Which normative distinctions must be kept separate when applying block?
\item How would you audit a program or implementation that depends on block?
\item Compare C standard library with character; what different question does each answer?
\item What role does character play in the clause's overall model?
\item Which normative distinctions must be kept separate when applying character container type?
\item How would you audit a program or implementation that depends on collating element?
\item Compare component with conditionally-supported; what different question does each answer?
\item What role does conditionally-supported play in the clause's overall model?
\item Which normative distinctions must be kept separate when applying constant subexpression?
\item Scenario: a program relies on an unstated implementation behavior while using deadlock. What must be checked before calling the program portable?
\item Scenario: a program relies on an unstated implementation behavior while using default behavior. What must be checked before calling the program portable?
\end{enumerate}
\clearpage\section{Answer Key}
\begin{enumerate}
\item access is one part of the clause's core-language semantics model. Its role is understood by connecting the local rule in 3.1 to the clause orientation and its dependent concepts.
\item Keep formation and well-formedness separate from runtime preconditions and effects. Also distinguish mandatory wording from notes, implementation choice from undefined behavior, and interface declarations from detailed contracts; see 3.5.
\item Audit argument by locating the controlling wording in 3.5, listing inputs and type constraints, proving any preconditions, recording effects and failure behavior, and checking lifetime, invalidation, complexity, and synchronization where applicable.
\item argument and argument occupy different steps or facility groups. Analyze each under its own subclause, then follow the explicit dependency between them rather than transferring one set of guarantees to the other; begin at 3.5.
\item argument is one part of the clause's core-language semantics model. Its role is understood by connecting the local rule in 3.5 to the clause orientation and its dependent concepts.
\item Keep formation and well-formedness separate from runtime preconditions and effects. Also distinguish mandatory wording from notes, implementation choice from undefined behavior, and interface declarations from detailed contracts; see 3.7.
\item Audit block by locating the controlling wording in 3.7, listing inputs and type constraints, proving any preconditions, recording effects and failure behavior, and checking lifetime, invalidation, complexity, and synchronization where applicable.
\item C standard library and character occupy different steps or facility groups. Analyze each under its own subclause, then follow the explicit dependency between them rather than transferring one set of guarantees to the other; begin at 3.8.
\item character is one part of the clause's core-language semantics model. Its role is understood by connecting the local rule in 3.9 to the clause orientation and its dependent concepts.
\item Keep formation and well-formedness separate from runtime preconditions and effects. Also distinguish mandatory wording from notes, implementation choice from undefined behavior, and interface declarations from detailed contracts; see 3.10.
\item Audit collating element by locating the controlling wording in 3.11, listing inputs and type constraints, proving any preconditions, recording effects and failure behavior, and checking lifetime, invalidation, complexity, and synchronization where applicable.
\item component and conditionally-supported occupy different steps or facility groups. Analyze each under its own subclause, then follow the explicit dependency between them rather than transferring one set of guarantees to the other; begin at 3.12.
\item conditionally-supported is one part of the clause's concurrency contract model. Its role is understood by connecting the local rule in 3.13 to the clause orientation and its dependent concepts.
\item Keep formation and well-formedness separate from runtime preconditions and effects. Also distinguish mandatory wording from notes, implementation choice from undefined behavior, and interface declarations from detailed contracts; see 3.14.
\item First identify the exact requirement in 3.15, then classify the relied-on behavior as required, unspecified, implementation-defined, conditionally supported, or undefined. Portability requires satisfying the program obligations and avoiding dependence on a choice not guaranteed in the target environments.
\item First identify the exact requirement in 3.16, then classify the relied-on behavior as required, unspecified, implementation-defined, conditionally supported, or undefined. Portability requires satisfying the program obligations and avoiding dependence on a choice not guaranteed in the target environments.
\end{enumerate}
\section{Key-Term Recap}
\begin{longtable}{@{}>{\RaggedRight\arraybackslash}p{0.28\textwidth}>{\RaggedRight\arraybackslash}p{0.64\textwidth}@{}}\toprule\textbf{Term} & \textbf{Working meaning}\\\midrule\endfirsthead\toprule\textbf{Term} & \textbf{Working meaning (continued)}\\\midrule\endhead\bottomrule\endfoot
well-formed program & A program satisfying the applicable syntax and semantic rules.\\
ill-formed program & A program that is not well-formed.\\
undefined behavior & Behavior for which no requirements are imposed.\\
unspecified behavior & A permitted implementation choice that need not be documented.\\
implementation-defined behavior & A permitted choice that each implementation documents.\\
valid but unspecified state & A state preserving invariants while leaving the value otherwise unspecified.\\
signature & The identity-relevant components of a function, member, template, or specialization in its stated context.\\
access & An execution-time read or modification of an object's value.\\
\end{longtable}
\CornellTakeaway{Precise vocabulary is part of the technical model: changing a defined term can change the conformance result.}
\end{document}
